\section{Introduction}\label{sec:intro}

Language-model agents now routinely act through tools whose effects cannot be undone---charging cards, moving money, sending messages, creating records~\cite{yao2023react,schick2023toolformer,yao2024taubench}. When such an agent crashes, times out, or is re-prompted after a failure, its recovery path re-issues work, and the surrounding system must not perform the same irreversible action twice. The classical defence is straightforward: give each effect an identity, record it when the effect commits, and refuse a second commit under the same identity~\cite{birrell1984rpc,kleppmann2017ddia,stripe2026idempotency}. Every mechanism of this kind---journal positions in durable-execution engines~\cite{temporal2026docs,restate2026docs}, client-chosen idempotency keys~\cite{stripe2026idempotency}, hashes of a canonicalised request---answers one question implicitly: \emph{when do two requests denote the same act?}

For a deterministic program the answer is positional or syntactic, because a retried call replays the original byte for byte. An LLM agent breaks both assumptions. Its control flow is sampled, so ``the third tool call'' is not stable across retries; and more fundamentally, its \emph{arguments} are re-synthesised. An agent asked to wire \$780 to Marlow Freight for shipment 88213 may first emit
\begin{quote}\small\ttfamily
wire\_money\{\\
\hspace*{1.5em}amount\_text: "\$780.00",\\
\hspace*{1.5em}beneficiary\_text: "Marlow Freight",\\
\hspace*{1.5em}reference: "Shipment 88213"\}
\end{quote}
and, after a crash and a re-prompt, emit
\begin{quote}\small\ttfamily
wire\_money\{\\
\hspace*{1.5em}amount\_text: "seven hundred eighty dollars",\\
\hspace*{1.5em}beneficiary\_text: "MARLOW FREIGHT",\\
\hspace*{1.5em}reference: "shipment \#88213"\}
\end{quote}
The two requests canonicalise differently and hash differently, and a content-addressed fence wires the money a second time. The converse situation is equally real: two requests that differ in a single token of free text (``instalment 1 of 2'' versus ``instalment 2 of 2'') can be two legitimate acts that must both execute.

Deciding identity therefore requires interpreting arguments, and the two errors are sharply asymmetric. Executing a duplicate is costly but \emph{observable}: two committed records exist for reconciliation to find. Suppressing a distinct act is \emph{silent}: the fence returns the recorded response of an earlier act, the agent believes the new act happened, and no trace shows otherwise. We call this second error \emph{wrong suppression} and treat it as the safety-critical direction throughout. Effect identity is therefore a problem of reasoning under uncertainty with asymmetric loss, not of bookkeeping.

\paragraph{The gap}
Canonical keys are sound but incomplete: under a semantics-preserving normalisation they never merge two distinct acts, yet they miss every paraphrased duplicate. Recent work identifies precisely this failure in agent checkpoint--restore (``semantic rollback attacks''~\cite{zheng2026acrfence}) and in authorisation (``semantic replay''~\cite{xu2026caplease}), and proposes comparing a re-synthesised call against logged calls with an analyser model~\cite{zheng2026acrfence}. An LLM judgement is the only general-purpose checker available for such predicates, but LLM judges are biased and promptable~\cite{zheng2023judge,wang2023fair,ye2024justice} and readily steered by adversarial text~\cite{greshake2023injection,hwang2025persuasion}. What is missing is a decision procedure whose safety-critical error is \emph{bounded} rather than inherited wholesale from an unconstrained judge.

\paragraph{Research question}
Can semantic effect identity be converted into a bounded decision mechanism---one that suppresses paraphrased duplicates while controlling the probability of wrongly suppressing a distinct act---and what does the declared tool schema determine about how far that is possible?

\paragraph{Approach}
We decompose effect identity into three parts and build a fence around the decomposition (Figures~\ref{fig:architecture} and~\ref{fig:flow}). The layers are consulted in order:
\begin{enumerate}[(1),leftmargin=*,itemsep=1pt]
\item a \emph{certifier}, a decidable relation that proves sameness (canonical equality under declared, semantics-preserving field policies); pairs it covers are suppressed without consulting any model;
\item a \emph{refuter}, a decidable relation that proves difference (inequality of declared, typed, effect-determining fields); pairs it covers are executed without consulting any model, so no model output---however manipulated---can suppress across a declared difference;
\item the \emph{residual}, the pairs neither settles; only these reach an LLM judge;
\item the judge operates under an asymmetric loss: it can cause suppression only when its score clears a threshold calibrated on distinct pairs, and scores in an intermediate band execute and are durably flagged;
\item the threshold is calibrated for the statistic the fence actually thresholds---the maximum judge score over the $m$ committed candidates that survive the refuter---with a one-sided split-conformal quantile at level $\alpha/m$;
\item the calibration pool is drawn from the rule's own jurisdiction (distinct pairs the rule would route to the judge), and when it is too small to certify level $\alpha$ at candidate-set size $m$ the threshold is $+\infty$, so the fence does not suppress semantically at all.
\end{enumerate}

\paragraph{Main theoretical results}
Theorem~\ref{thm:bound} establishes that if the judge scores of truly distinct residual comparisons are exchangeable with the calibration pool, conditionally on $m$, then for a request denoting a different act from every one of its $m$ surviving candidates the probability of suppression is at most $\alpha$---with no assumption on the score distribution and no assumption on the dependence among the $m$ comparisons. The guarantee is marginal over calibration draws. It does not assert that the judge errs less than $\alpha$ of the time, nor that the rate realised by any particular deployment is below $\alpha$. Proposition~\ref{prop:residual} shows that the residual is a function of the declared schema alone: no judge moves a pair out of it, and every error of a mechanism with sound decidable layers lies inside it. Proposition~\ref{prop:dilemma} gives the resulting \emph{refinement dilemma} in closed form: refining the schema shrinks the residual and therefore the calibration pool, moving the conformal threshold to a higher order statistic and shrinking the certifiable candidate budget $\lfloor\alpha(n+1)\rfloor$. Schema refinement buys decidable safety at a cost the statistical layer pays.

\paragraph{Main empirical results}
On 1,066 constructed pairs over six irreversible tools, evaluated with three LLM judges from two model families (Section~\ref{sec:results}):
\begin{itemize}[leftmargin=*,itemsep=1pt]
\item the fence deduplicates 94.7\%~[91.4--96.7] of paraphrased duplicates at 1.7\%~[0.7--4.3] wrong suppression with the primary judge at $\alpha=0.05$, against 12.1\%~[8.8--16.4] for canonical keys;
\item wrong suppression stays below $\alpha$ for all three judges (1.7\%, 1.3\%, 0.9\%) while deduplication ranges from 82.6\% to 94.7\%, and the three judges' wrong suppressions fall on disjoint pairs;
\item in a controlled simulation, thresholding the maximum of $m=50$ candidate scores with an uncorrected pairwise threshold wrongly suppresses 89.0\% of distinct requests, against 0.0\% for the corrected rule (which switches off beyond its budget);
\item the refuter removes a typed-field injection attack that the raw primary judge accepts 65.0\% of the time, but prose-only injection still succeeds 84.0\% of the time;
\item refining the declarations so that the numbers mentioned in free-text fields become typed closes the prose-only injection class, but lowers deduplication to 69.0\% on the primary judge and to 12.1\% on the other two; the smaller calibration pool and its changed composition both contribute.
\end{itemize}

\paragraph{Important limitations}
The labels operationalise the authors' definition of identity. An independent human annotation was designed, preregistered and instrumented but not collected; a blinded panel of language models was run instead and is reported only as model-panel evidence, alongside two defects in its preregistered re-derivation discovered during this revision (Section~\ref{sec:res-labels}). Under seed-level resplitting, mean wrong suppression rises from 1.7\% to 2.7\% on the primary judge, and 10.5\% of splits exceed $\alpha$. Distinctions carried only by free text remain judge-dependent. The mechanism provides bounded identity recognition and fence-level at-most-once commitment under stated assumptions; it does not provide exactly-once execution at an external provider (Section~\ref{sec:guarantees}).

\begin{center}
\fbox{\parbox{0.92\linewidth}{\textbf{Scope.} The proposed mechanism is a bounded semantic identity and duplicate-suppression layer. It is not a universal exactly-once execution protocol.}}
\end{center}

\paragraph{Contributions}
\begin{enumerate}[(C1),leftmargin=*,itemsep=1pt]
\item A formalisation of semantic effect identity for irreversible agent actions under an asymmetric loss, which separates identity recognition, fence-level at-most-once commitment and external exactly-once execution (Section~\ref{sec:problem}).
\item The certifier/refuter/residual decomposition and the observation that the residual is determined by the declared schema (Section~\ref{sec:framework}).
\item A fence that consults an LLM judge only on the residual and calibrates the deployed max-over-candidates statistic with a multiplicity-corrected one-sided conformal threshold on a jurisdiction-matched pool, with a finite-sample bound on wrong suppression and an explicit calibration budget (Sections~\ref{sec:fence} and~\ref{sec:theory}).
\item The refinement dilemma: an exact statement of how schema refinement shrinks the calibration budget (Proposition~\ref{prop:dilemma}) and an empirical decomposition of its pool-size and pool-composition components (Section~\ref{sec:res-refinement}).
\item An evaluation designed around the framework's boundaries---three judges, four paraphrase-generator arms including a hand-written one, adversarial injection, schema refinement, crash and replay across real process death, relabelling and resplitting---with an artifact whose build fails when a quoted interval does not match a computed one (Sections~\ref{sec:method} and~\ref{sec:results}).
\end{enumerate}
None of the individual techniques---canonicalisation, typed validation, LLM judging, split-conformal calibration, Bonferroni correction---is new. The contribution lies in their combination into a decision procedure for effect identity and in the analysis of what that combination can and cannot guarantee. The underlying pattern---an autonomous agent making a control decision with asymmetric consequences, a stochastic oracle as the only general-purpose decider, and adversarial input---recurs across agent safety, and the recipe may transfer (Section~\ref{sec:discussion}).

Section~\ref{sec:related} positions the work. Section~\ref{sec:problem} formalises the problem and the three guarantees, Section~\ref{sec:framework} the decomposition, Section~\ref{sec:fence} the fence, and Section~\ref{sec:theory} its analysis. Sections~\ref{sec:method} and~\ref{sec:results} report the evaluation, Section~\ref{sec:discussion} discusses it, and Section~\ref{sec:limitations} states limitations and broader impact. Proofs are in \ref{app:proofs}; secondary experiments, full tables and reproduction details are in the supplementary material.
